%%% Colour commands - delete for final draft
% --- Editing Macros (Hierarchy of Importance) ---
% 1. FACT CHECK (Critical): Red commands immediate attention for potential inaccuracies.
\newcommand{\fc}[1]{\textcolor{red}{[FactC: #1]}}
% 2. CITE (High): Orange highlights missing evidence or attribution.
\newcommand{\ct}[1]{\textcolor{orange}{[CITE: #1]}}
% 3. APPENDIX (Medium): Blue indicates structural moves (shifting text out of the main body).
\newcommand{\app}[1]{\textcolor{blue}{[APP: #1]}}
% 4. SMALL REWORD (Low): Teal is for minor stylistic edits, mirroring your friend's setup.
\newcommand{\sr}[1]{\textcolor{teal}{[Sreword: #1]}}
% 5. FIGURE (Low): Purple is for figure-related notes (refer to figure, create figure, etc.).
\newcommand{\fig}[1]{\textcolor{purple}{[FIG: #1]}}
% 6. AI redraft (Low): Ai marker suggestions and corrections
\newcommand{\ai}[1]{\textcolor{brown}{[AI: #1]}}


%=======================================================================================
\chapter{Related Work and Background}
\label{ch:background}

This chapter reviews the literature in three areas: top-height and growth-curve modelling, machine learning and physics-informed prediction, and spatial evaluation and attribution methods. It closes by identifying where this literature falls short and how the study addresses that gap.


%------------------------------------------------
\section{Top-height growth and environmental variation}
\label{sec:background-top-height}

\paragraph*{Chapman--Richards growth models.}
The Chapman--Richards (CR) equation is widely adopted in forestry to represent height and yield trajectories through a flexible sigmoidal formulation \citep{richards1959Flexible_C2_CR,pommerening2015Methods_C2_CR,manso2021Dynamic_C2_CR}. In this context, the model is parameterised as:
\begin{equation}
H(a) = y_{\max}\left(1-\exp(-ka)\right)^p,
\label{eq:background-cr}
\end{equation}
where $H(a)$ is dominant height at stand age $a$, $y_{\max}$ is the asymptotic height ceiling, $k$ controls the approach rate to the asymptote, and $p$ governs the inflection and shape of the trajectory.

% \fig{FIGURE: Plot one illustrative Chapman--Richards curve (Equation~\ref{eq:background-cr}), height (m) on the y-axis, stand age (years, 0--100) on the x-axis. Use plausible parameter values only; state directly in the caption that this is a representative example curve, not a fit to this project's own data, and that the exact values do not need justifying. Draw the curve as a single solid line in one clear colour (e.g.\ dark green). Annotate the asymptote $y_{\max}$ with a horizontal dashed line, labelled to state it represents the stand's maximum height potential under the fitted curve, not a height any individual stand is necessarily observed at. Add a second horizontal reference line in a visually distinct colour (e.g.\ amber), labelled with Sitka spruce's typical UK top height at commercial felling age: 16--23\,m at yield class 14, around 40 years, citing the same Forest Research yield-table source already used elsewhere in this dissertation \cite{forestResearchForestYield}. 60\,m as  local ecological maximum for Sitka spruce in Aberfoyle reference \cite{mason2011Sitka_C1_SS}; if a maximum-height marker is added, use that figure and citation, not a new uncited number. The caption must state: (1) this is an illustrative curve, not a fitted result; (2) what the asymptote represents; (3) that stands are typically felled well before the curve approaches its asymptote, which is why the asymptote itself, not any single observed height, is the growth-potential parameter this dissertation models; (4) the felling-height source and figure.}

When observational data span a restricted age interval, the asymptote $y_{\max}$ becomes difficult to isolate, allowing different combinations of $(y_{\max}, k, p)$ to produce overlapping trajectories. Factors such as measurement noise, repeated-survey correlation, and stand disturbance cannot be resolved simply by choosing a flexible curve, so a single shared curve represents the average trajectory of the entire population \citep{rennolls1995Forest_C2_CR}.

% \ai{Equation~\ref{eq:background-cr}, the parameter definitions and the general strengths and limitations of Chapman--Richards can be removed from Methodology. That chapter should retain only the training-partition fitting procedure, parameter restrictions and the way the fitted relationship enters each experiment.}

\paragraph*{Top height and forest productivity.}
British forestry management decisions draw heavily on this same height--age relationship, for estimations such as General Yield Class (GYC), a measure of site productivity \citep{edwards1981Yield_C1_SS,ForestResearch2026Sitka_C1_SS}. GYC defines the maximum expected mean annual increment in cumulative stemwood timber volume ($\mathrm{m^3\,ha^{-1}\,yr^{-1}}$), scaling top height by the same Chapman--Richards expansion factor introduced above:
\begin{equation}
\mathrm{GYC}_{\mathrm{SS}} = \frac{\mathrm{TopHeight}_{\mathrm{SS}} \left(1 - \exp(-0.035067\,\mathrm{Age}_{\mathrm{SS}})\right)^{-1.923684} - 12.13610}{1.473882}
\label{eq:gyc-ss}
\end{equation}
where $\mathrm{TopHeight}_{\mathrm{SS}}$ is dominant canopy height (m) and $\mathrm{Age}_{\mathrm{SS}}$ is stand age (years). Stands reaching a greater top height at a given age receive a higher expected productivity rating, which informs downstream forecasts of merchantable timber volume.

Using inventory metrics like GYC or timber volume to predict top height causes target leakage, as both are derived from height via \eqref{eq:gyc-ss}. Furthermore, while standard yield curves \citep{edwards1981Yield_C1_SS} offer reliable regional baselines, they show average trends and miss plot-to-plot differences.

% \ai{The general forestry justification for selecting top height can be removed from Methodology. The Data chapter should retain only the exact definition of the LiDAR-derived response and its relationship with the original Forest Research fields. - also the motivation is mentioned in introduction, so things in this chapter should focus more on the papers and documents etc and final conclusions of reasons don't have to be mentioned until research gap}

\paragraph*{Regional and environmental variation in height growth.}
Stands with similar early growth often diverge later in life due to differences in local environment \citep{socha2021Regional_C2_CR}. Site-conditioned and mixed-effect models account for this split by adjusting growth-curve parameters using measured plot or stand variables \citep{scolforo2013Dominant_C2_CR,socha2016Assessment_C2_CR}.

For Sitka spruce in the UK and Ireland, site productivity depends heavily on soil moisture, elevation, wind exposure, and nutrient availability \citep{worrell1990Productivity1_C1_SS,farrelly2011Sitka_C1_SS}. Upland surveys show that General Yield Class drops by $3.2$--$4.0\,\mathrm{m^3\,ha^{-1}\,yr^{-1}}$ for every 100\,m rise in elevation, mostly tracking shifts in temperature and wind. Across Irish plantations, soil nutrient regime accounts for over half the variation in site index, while growth slows at both dry and wet moisture extremes, and topographic exposure becomes limiting only on high-elevation sites \citep{farrelly2011Sitka_C1_SS}. Because these factors correlate, single coefficients often reflect combined site conditions, not isolated effects. Unlike direct field surveys, coarse digital rasters risk masking local environmental signals; for instance, global SoilGrids layers correlate weakly with ground measurements ($R^2 = 0.19$; \citealt{miller2024Assessing_C2_CR}), although local high-resolution rasters retain stronger value.

\paragraph*{Multi-source environmental modelling.}
Integrating forest inventory records with multi-source geospatial data—such as digital elevation models, soil classifications, and topographic exposure indices—provides an established framework for spatial productivity mapping. Early UK applications linked Forestry Commission subcompartment databases with terrain and soil layers via principal component analysis and global linear regression to map Sitka spruce Yield Class across regional scales \citep{bateman1998Using_C1_SS}, though this evaluated coarse stand-level productivity rather than repeated plot-level height trajectories.

More recent remote-sensing workflows combine airborne LiDAR canopy height metrics with multi-scale climate, soil, and topographic layers. In complex terrain, macro-climatic indices (e.g., climatic water balance) can govern broad-scale maximum canopy height, whereas fine-scale topographic variables explain localized structural variation \citep{fricker2019More_C3_LID}. This confirms that the apparent explanatory power of environmental covariates is highly sensitive to the spatial resolution and alignment of both the response metric and the predictor layers.

% \ai{The literature-based argument that growth relationships vary with site conditions can be removed from the environmental-model Methodology. Methodology should retain only which environmental features were supplied, which Chapman--Richards parameters were conditioned and how the adjustment was implemented.}


%------------------------------------------------
\section{Data-driven and process-informed prediction}
\label{sec:background-prediction}

Machine learning models learn patterns directly from data without assuming a fixed curve shape \citep{lee2018Machine_C3_LID}. Tree ensembles are standard, effective baselines for structured tabular data. Random Forest builds an ensemble of independent decision trees through bagging and random feature selection, averaging their outputs to reduce variance \citep{breiman2001RandomForest_C4_NN}. In contrast, Gradient Boosted Decision Trees, such as XGBoost, build trees sequentially where each new tree corrects the residual errors of the previous ones \citep{chen2016XGBoost_C4_NN}. Both approaches handle non-linear relationships and feature interactions well without requiring manual parameter curves. However, forest remote-sensing studies often test these models on single-date snapshots; predicting repeated height growth requires models to generalize to unseen forest blocks, across stand ages, and over multi-year survey intervals.

\paragraph*{Neural models of forest height growth.}
Deep neural networks (DNNs) learn unconstrained multi-variable interactions via layered hierarchical representations, optimising their weights through backpropagation to minimise a specified training loss. However, unconstrained empirical networks do not inherently preserve theoretical growth properties, such as monotonic height accumulation, bounded asymptotes, or base-age invariance \citep{ercanli2023Comparison_C4_NN,karatepe2022Total_C4_NN}. When evaluated against non-linear mixed-effects models, unconstrained ANNs can reproduce general sigmoid shapes while failing to preserve essential demographic properties across age ranges \citep{ercanli2023Comparison_C4_NN}.

% \ai{Generic neural-network theory and the general justification for a DNN can be removed from Methodology. That chapter should retain the actual network architecture, inputs, loss and training procedure.}

\paragraph*{Process-informed neural models.}
Physics-informed neural networks (PINNs) and theory-guided machine learning frameworks incorporate scientific relationships directly into the optimization objective \citep{karpatne2017Theory_C4_NN,raissi2019Physics_C4_NN,karniadakis2021Physics_C4_NN,willard2022Integrating_C4_NN}. A standard loss-regularised formulation is expressed as:
\begin{equation}
\mathcal{L}_{\text{total}} = \mathcal{L}_{\text{data}} + \lambda\,\mathcal{L}_{\text{process}},
\label{eq:pinn-loss-general}
\end{equation}
where $\mathcal{L}_{\text{data}}$ measures observational error, $\mathcal{L}_{\text{process}}$ penalizes departures from domain-specific equations, and $\lambda$ balances the objectives.

Process knowledge can be integrated into neural architectures across a structural continuum \citep{wesselkamp2024Process_C4_NN}:
\begin{table}[ht]
\centering
\scriptsize
\begin{tabular}{ll}
\hline
\textbf{Integration Strategy} & \textbf{Structural Mechanism} \\
\hline
\emph{Bias correction} & Network learns systematic residuals of process model outputs. \\
\emph{Parallel physics} & Process and neural predictions are combined additively. \\
\emph{Physics regularisation} & Unconstrained neural predictor is penalised for equation disagreement. \\
\emph{Domain adaptation} & Pre-training on process simulations before empirical fine-tuning. \\
\emph{Physics embedding} & Differentiable process models embedded with parameter sub-networks. \\
\hline
\end{tabular}
\label{tab:pinn-strategies}
\end{table}
Conditional and variable-coefficient PINN frameworks combine physics regularisation with physics embedding, letting environmental covariates drive growth-curve parameters rather than treating them as fixed constants. Frameworks such as Forest-Informed Neural Networks integrate neural networks into more complex differential equation models covering growth, mortality, regeneration, and competition; a single Chapman--Richards curve is a simpler target, for which the ecological equation can instead be anchored with a soft penalty term \citep{pichler2025Inferring_C4_NN}.

Adding process penalties creates a trade-off between fitting the observations and adhering to the Chapman--Richards guidance. Composite loss functions can produce imbalanced gradients, causing one loss component to dominate optimisation, while ill-conditioned loss landscapes can slow or prevent convergence \citep{wang2021Understanding_C4_NN,rathore2024Challenges_C4_NN}. There is also a risk that a strong penalty based on a misspecified population-level Chapman--Richards curve pulls predictions towards an unsuitable average rather than allowing the network to represent local growth differences.

%------------------------------------------------
\section{Spatial evaluation and growth attribution}
\label{sec:background-spatial-attribution}

Evaluating spatial models requires separating temporal forecasting from geographical transfer. Nearby inventory plots share climate, soil, planting history, and remote-sensing capture conditions, causing strong spatial autocorrelation. While holding out future surveys tests temporal extrapolation within a known forest, evaluating on disjoint spatial blocks reveals whether the model can generalise across regional gradients \citep{ploton2020Spatial_C5_EVAL}.

Predictive accuracy and environmental attribution are different goals. A variable can improve height predictions by tracking broad trends without reflecting the underlying biological driver. Attribution aims to uncover cause-and-effect relationships, either through transparent models with directly readable coefficients, such as a regularised linear benchmark combining $L_1$ and $L_2$ penalties for automated feature selection, or by applying post hoc interpretation tools like SHAP to more flexible models such as gradient-boosted trees \citep{lundberg2017Unified_C7_XAI}. Grounded in cooperative game theory, SHAP calculates the marginal contribution of each predictor to the final prediction by evaluating model outputs across all possible feature combinations. But these attribution tools can only quantify statistical associations with model outputs, not causes, and this limitation is amplified by spatial mismatch: external gridded climate and soil data are often too coarse to capture plot-scale conditions, missing localised competition or wind protection from nearby trees. A regularised linear model additionally assumes the environmental relationship is identical everywhere in the study area, an assumption spatially varying models relax further below.

Local regression models relax this assumption, allowing coefficients to vary by location instead of taking one fixed value for the whole study area. Where standard Geographically Weighted Regression (GWR) relies on fixed distance-decay kernels and bandwidths \citep{brunsdon1996Geographically_C6_GWR}, GNNWR uses a neural network to learn dynamic spatial weighting functions \citep{du2020Geographically_C6_GWR}. This preserves transparent local parameters while capturing complex regional variation, albeit at the cost of high pairwise distance computation and memory demand on large datasets.

%------------------------------------------------
\section{Research gap and study position}
\label{sec:background-gap}

Established top-height models provide simple representations of stand development, yet uniform height-age curves often mask significant local growth variation. Multi-temporal LiDAR data provides extensive spatial coverage of dominant canopy height, but these metrics reflect a combination of biological growth, sensor acquisition differences, and unmeasured plot-level management.

ML has already been used to relate LiDAR-derived canopy height to environmental conditions in other forests \citep{fricker2019More_C3_LID}. However, no study identified in this review combines repeated LiDAR observations with stand records and GIS-derived terrain, wind, climate and soil variables to model Sitka spruce height trajectories at Aberfoyle, where direct field measurements are limited. ML and PINNs provide flexible alternatives to static parametric curves, but adding domain guidance does not inherently guarantee better predictions than an unconstrained DNN or a tuned tree ensemble. Separately, the attribution of plot-level growth-curve departures has not been compared across linear, tree-based and spatially varying models under spatially independent validation, leaving it uncertain whether different model families identify consistent environmental influences.

This dissertation addresses these gaps through the aim and objectives set out in the introduction.


% Explain what previous studies did.
% State the important findings.
% Compare the approaches.
% Identify limitations that lead directly to the dissertation’s research gap.